\section{Resultados experimentales y conclusiones prácticas}

\subsection{Interpretación de los resultados: la tendencia importa más que los valores absolutos}
El curso muestra una tendencia de mejora en varios benchmarks, especialmente en tareas más difíciles y en configuraciones en línea, donde agregar planificación y una gobernanza más sólida de los datos de entrenamiento hace que el efecto sea más estable. Aunque las puntuaciones concretas varían según la configuración, la tendencia se mantiene muy consistente.

\begin{figure}[H]
\centering
\includegraphics[width=0.88\textwidth]{frame-08-results.jpg}
\caption{Video aproximadamente en 01:24:55: la página de resultados destaca cómo el entrenamiento por etapas y el planner elevan la tasa de éxito en tareas complejas}
\end{figure}

\begin{importantbox}{Tres lecciones de ingeniería reutilizables}
\begin{enumerate}
    \item Primero hay que estabilizar la evaluación y el entorno, y solo después hablar de la velocidad de iteración del modelo;
    \item la prioridad de la gobernanza de datos no es menor que la de las mejoras del modelo;
    \item para tareas complejas, la señal de planificación suele ser una ``condición necesaria'' y no ``un simple adorno''.
\end{enumerate}
\end{importantbox}

\subsection{Recomendaciones de métricas para la puesta en producción}
\begin{table}[H]
\centering
\begin{tabular}{p{3.2cm}p{5.0cm}p{2.8cm}}
\toprule
\textbf{Categoría de métrica} & \textbf{Métrica sugerida} & \textbf{Escenario de uso} \\
\midrule
Resultado de la tarea & Tasa de finalización, tasa de finalización parcial, distribución de tipos de fallo & Evaluar la utilidad del producto \\
Calidad de la trayectoria & Número promedio de pasos, tasa de acciones repetidas, tasa de retrocesos & Diagnosticar la calidad de la planificación \\
Estabilidad del sistema & Consistencia entre reejecuciones, tasa de recuperación tras interrupciones anómalas & Evaluar la confiabilidad en producción \\
Ciclo cerrado de datos & Proporción de muestras fallidas reincorporadas, magnitud de la mejora tras la corrección & Medir la eficiencia de la iteración \\
\bottomrule
\end{tabular}
\caption{Marco de métricas de ingeniería de agentes derivado de las conclusiones del curso}
\end{table}

\begin{warningbox}{No te fijes solo en la ``tasa de finalización total''}
Una sola métrica puede ocultar problemas. Por ejemplo, un aumento en la tasa de finalización puede deberse a que aumentó la proporción de tareas simples, y no a una mejora en la capacidad para tareas complejas. El curso recomienda hacer al menos un desglose estadístico por dificultad de la tarea, tipo de tarea y tipo de entorno.
\end{warningbox}

\subsection{Resumen de la sección}
Lo más valioso de los resultados de esta clase es la metodología: el entorno, la evaluación, los datos y el modelo deben avanzar de manera simultánea. Cualquier eslabón débil hará que el sistema se vuelva inestable en tareas reales.
